\documentclass[a4paper]{article}

%------------------------------------------------
% الحزم
%------------------------------------------------

\usepackage{enumitem}
\usepackage[margin=2.5cm]{geometry}
\usepackage{amsmath,amsfonts,amssymb}
\usepackage[charter]{mathdesign}
\usepackage{graphicx}
\usepackage{array}
\usepackage{fancyhdr}
\usepackage{polyglossia}

%------------------------------------------------
% اللغات والخطوط
%------------------------------------------------

\setdefaultlanguage[
    calendar=gregorian,
    numerals=maghrib
]{arabic}

\setotherlanguage{english}

% إذا لم تكن الخطوط مثبتة في الجهاز، تُؤخذ من المجلد ../fonts/
\IfFontExistsTF{Amiri}{%
    \newfontfamily\arabicfont[Script=Arabic, Scale=1.15]{Amiri}
}{%
    \newfontfamily\arabicfont[
        Script=Arabic, Scale=1.15,
        Path=../fonts/, Extension=.ttf,
        UprightFont=*-Regular, BoldFont=*-Bold,
        ItalicFont=*-Italic, BoldItalicFont=*-BoldItalic
    ]{Amiri}
}

\IfFontExistsTF{AlBattar}{%
    \newfontfamily\arabicfontsf[Script=Arabic, Scale=1.2]{AlBattar}
    % خط العناوين (يعمل حتى بعد استعمال \textenglish)
    \newfontfamily\khatt[Script=Arabic, Scale=1.2]{AlBattar}
}{%
    \newfontfamily\arabicfontsf[
        Script=Arabic, Scale=1.2, Path=../fonts/,
        BoldFont=ae_AlBattar.ttf
    ]{ae_AlBattar.ttf}
    % خط العناوين (يعمل حتى بعد استعمال \textenglish)
    \newfontfamily\khatt[
        Script=Arabic, Scale=1.2, Path=../fonts/,
        BoldFont=ae_AlBattar.ttf
    ]{ae_AlBattar.ttf}
}

%------------------------------------------------
% إعداد الصفحة
%------------------------------------------------

\renewcommand{\footrulewidth}{1pt}
\renewcommand{\headrulewidth}{1pt}

\rhead{\khatt المدرسة العليا للأساتذة ببوزريعة}
\lhead{\khatt ملحقة بومرداس}

\cfoot{}
\lfoot{\khatt الأستاذ: د. بن عبد الرزاق يوسف}

\parindent 0pt
\setlength{\headsep}{15pt}
\setlength{\parskip}{3pt}

\pagestyle{fancy}

\setlist[enumerate,1]{
    label=\arabic*),
    itemsep=4pt,
    topsep=3pt
}

\setlist[itemize]{
    itemsep=2pt,
    topsep=2pt
}

\arraycolsep=3pt

%================================================
% بداية الوثيقة
%================================================

\begin{document}

%------------------------------------------------
% العنوان
%------------------------------------------------

\begin{center}

{\Large\khatt\textbf{سلسلة تمارين رقم 3}}

\vspace{0.15cm}

{\Large\khatt\textbf{التفاعل الفعال مع الذكاء الاصطناعي}}

\vspace{0.1cm}

{\khatt\textbf{مبادئ هندسة الأوامر ــ Prompt Engineering}}

\end{center}

\rule{\textwidth}{1pt}

\begin{center}

{\khatt\textbf{
الفصل 3: التفاعل الفعال مع الذكاء الاصطناعي
}}

\end{center}

\rule{\textwidth}{1pt}

%================================================
% وضعية الانطلاق
%================================================

\section*{\khatt\underline{وضعية محفزة }}

كتب الطالب الأول:



\textbf{
``أعطني خطة للمراجعة.''
}


بينما كتب الطالب الثاني:

\begin{center}

\textbf{
``تصرف كمساعد أكاديمي لطالب جامعي مبتدئ.
أنشئ خطة مراجعة لمادة الذكاء الاصطناعي لمدة 7 أيام،
بمعدل ساعتين يوميًا، مع تحديد موضوع كل يوم،
وثلاثة أهداف تعليمية، وتمرين قصير في نهاية كل يوم.
نظم الإجابة في جدول واضح.''
}

\end{center}

\textbf{السؤال:}
لماذا يمكن أن تختلف جودة الإجابتين رغم أن الموضوع نفسه؟

\vspace{0.2cm}

%================================================
% النشاط الأول
%================================================

\section*{\khatt\underline{النشاط الأول: من الأمر العشوائي إلى الأمر الهندسي}}

\begin{enumerate}

%------------------------------------------------

\item جرّب الطلب التالي باستعمال أداة للذكاء الاصطناعي:

\begin{center}

\textbf{
``اشرح الذكاء الاصطناعي.''
}

\end{center}

سجّل الإجابة التي حصلت عليها.

%------------------------------------------------

\item أعد صياغة الطلب بإضافة \textbf{السياق}:

\begin{center}

\textbf{
``أنا طالب جامعي مبتدئ في الذكاء الاصطناعي.
اشرح لي مفهوم الذكاء الاصطناعي بطريقة بسيطة.''
}

\end{center}

قارن بين الإجابتين.

%------------------------------------------------

\item طوّر الطلب السابق بإضافة \textbf{المهمة} و
\textbf{القيود}:

\begin{center}

\textbf{
``أنا طالب جامعي مبتدئ في الذكاء الاصطناعي.
اشرح لي مفهوم الذكاء الاصطناعي في 150 كلمة،
بلغة عربية بسيطة، مع مثالين من الحياة اليومية.''
}

\end{center}

ما التغييرات التي لاحظتها؟

%------------------------------------------------

\item أضف إلى الطلب السابق \textbf{تنسيق الإجابة}:

\begin{center}

\textbf{
``نظم الإجابة في ثلاثة أجزاء:
تعريف، أمثلة، وخلاصة قصيرة.''
}

\end{center}

هل أصبحت الإجابة أكثر تنظيمًا؟ علّل.

%------------------------------------------------

\item أكمل الجدول التالي:

\begin{center}

\begin{tabular}{|
>{\raggedleft\arraybackslash}m{4cm}|
>{\raggedleft\arraybackslash}m{10cm}|}

\hline

\textbf{العنصر} & \textbf{دوره في الأمر} \\

\hline

المهمة
&
ما الذي نريد من الذكاء الاصطناعي أن يقوم به؟
\\
\hline

السياق
&
ما المعلومات التي تساعد الأداة على فهم الوضعية؟
\\
\hline

الدور
&
ما الشخصية أو الخبرة التي نطلب من الأداة اعتمادها؟
\\
\hline

القيود
&
ما الشروط التي يجب احترامها أثناء الإجابة؟
\\
\hline

التنسيق
&
كيف نريد أن تكون الإجابة منظمة أو معروضة؟
\\
\hline

\end{tabular}

\end{center}

%------------------------------------------------

\item استخرج من الطلب التالي عناصره الأساسية:

\begin{center}

\textbf{
``تصرف كأستاذ جامعي في البرمجة.
اشرح لطالب مبتدئ مفهوم الحلقة التكرارية
باستعمال لغة C، في أقل من 200 كلمة،
مع مثال بسيط وجدول يوضح خطوات التنفيذ.''
}

\end{center}

حدّد:

\begin{itemize}

    \item الدور:
    \item المهمة:
    \item السياق:
    \item القيود:
    \item التنسيق المطلوب:

\end{itemize}

\end{enumerate}

%================================================
% النشاط الثاني
%================================================

\section*{\khatt\underline{النشاط الثاني (واجب منزلي): تحديد الدور والتنسيق}}

\textbf{يُنجز هذا النشاط في المنزل ويُسلَّم في الحصة القادمة.}

\begin{enumerate}

\item اختر موضوعًا أكاديميًا تعرفه،

اكتب أمرًا يطلب من الذكاء الاصطناعي أن يتصرف كخبير
أو أستاذ في هذا المجال.

%------------------------------------------------

\item أنشئ أمرًا آخر يحدد الجمهور المستهدف:

\begin{center}

\textbf{
``اشرح الموضوع لطالب جامعي مبتدئ...''
}

\end{center}

%------------------------------------------------

\item اطلب من الأداة إنتاج الإجابة في أحد التنسيقات التالية:

\begin{itemize}

    \item جدول.
    \item نقاط مختصرة.
    \item خطوات مرقمة.
    \item مثال تطبيقي.

\end{itemize}

%------------------------------------------------

\item قارن النتائج.

ما أثر تحديد \textbf{الدور} و\textbf{الجمهور} و
\textbf{التنسيق} في جودة المخرجات؟

\end{enumerate}


\end{document}